% Figures du chapitre « Le génome synonyme, mode rapide » du Tome III (lot B de la construction,
% genome/07-genome-synonyme.md ; juge : genome/JUGEMENT-construction.md). Gardé par \ifdefined\GrapheSynoBasePropre.
% Données : \Racine/figures/tome3/donnees/synonyme-*.csv et synonyme-nombres.tex, écrits par
% memoire/outils/tome3_synonyme.py depuis mesures/genome/construction/{synonyme,juge}/ et mesures/moteur/O/ (C4).
% Styles préfixés « syno ».
%
% Second flux (2026-10-27) : couleurs, traits et marques du dictionnaire du Tome III (palette-tome3.tex : frontière,
% synonymes λ = 1 et λ = 0, Q2_0, moteur exact, R1 et fenêtres neuves, hors domaine) ; légendes G9 (\Titre,
% \Legende, \Source ; le chapitre écrit \LegendeFigure{<Nom>}). Graves corrigés : G3 (l'aire n'est pas le
% sur-apprentissage), G4 (la queue repose sur trois positions), G1 et G5 (III-9.3 : un seul axe, aucun point hors
% du cadre, chemins mêlés dits), G3, G4 et G7 du juge de design.
%   III-9.1 (a)(b)(b′) SynoBasePropre   III-9.2 SynoQueue   III-9.3 SynoForet   III-9.6 SynoVariantes
% Le panneau (c) de III-9.1 est la rangée de 32 poids du lot « espaces latents » : ce fichier nomme sa part
% cumulée (b′) pour lui laisser la lettre (c).

\ifdefined\GrapheSynoBasePropre\else
\input{\Racine/figures/tome3/legende-tome3}
\providecommand{\CodonDonnees}{\Racine/figures/tome3/donnees}
\providecommand{\CodonNum}[1]{\num[group-digits=integer]{#1}}
\input{\CodonDonnees/synonyme-nombres}

\tikzset{
  % Calibration contre évaluation (III-9.1) : le trait plein est la calibration, le pointillé ce que voit
  % l'évaluation ; la couleur reste celle de la variante.
  syno calib/.style={solid, line width=1.1pt},
  syno eval/.style={dash pattern=on 1.2pt off 1.6pt, line width=1.3pt},
}

% =============================================================================================================
% III-9.1 (a)(b)(b′). La mutation synonyme dans la base propre de H (L8, E0, gate du 20b).
\newcommand{\SynoMuet}[2]{% les directions muettes en calibration (k ≥ rang), hors du domaine de H : hachurées
  \fill[t3 hors domaine] (axis cs:\SynoNRang,#1) rectangle (axis cs:2880,#2);}
\newcommand{\GrapheSynoBasePropre}{%
\begin{tikzpicture}
% Troisième passe (DN16 du juge de design) : 4 mm de plus entre (a) et (b), pris sur leurs largeurs.
\begin{axis}[t3 graphe, name=a, width=0.485\linewidth, height=5.4cm, ymode=log,
  title={(a) \EF{the spectrum of $H$}{le spectre de $H$}}, xtick={0,500,...,2500},
  xlabel={\EF{eigendirection $k$ (by decreasing $\lambda_k$)}{direction propre $k$ (par $\lambda_k$ décroissant)}},
  ylabel={\EF{weight of direction $k$}{poids de la direction $k$}},
  xmin=0, xmax=2880, ymin=1e-3, ymax=1e3, unbounded coords=jump, legend pos=north east]
\SynoMuet{1e-3}{1e3}
\addplot[Encre, syno calib] table[x=k, y=lambda_calib] {\CodonDonnees/synonyme-spectre.csv};
\addlegendentry{$\lambda_k$ (\EF{calibration}{calibration})}
\addplot[Encre!70, syno eval] table[x=k, y=lambda_eval] {\CodonDonnees/synonyme-spectre.csv};
\addlegendentry{$u_k^\top H_{\EF{\mathrm{ev}}{\mathrm{év}}}\,u_k$ (\EF{evaluation}{évaluation})}
\draw[black!55, line width=0.5pt] (axis cs:\SynoNRang,1e-3) -- (axis cs:\SynoNRang,1e3);
\node[t3 note, anchor=south west] at (axis cs:\SynoNRang+40,1.6e-3)
  {\EF{rank}{rang} \CodonNum{\SynoNRang}\\\EF{silent in calibration}{muettes en calibration}};
\end{axis}
\begin{axis}[t3 graphe, name=b, at={(a.east)}, anchor=west, xshift=1.9cm, width=0.485\linewidth, height=5.4cm, ymode=log,
  title={(b) \EF{where the weight error goes}{où va l'erreur sur les poids}}, xtick={0,500,...,2500},
  xlabel={\EF{eigendirection $k$}{direction propre $k$}},
  ylabel={$e_k=\lVert\Delta W u_k\rVert^2$ (\EF{mean over}{moyenne par} \SynoNPaquet)},
  xmin=0, xmax=2880, ymin=1e-3, ymax=0.1, legend pos=south east]
\SynoMuet{1e-3}{0.1}
\addplot[t3 frontiere, t3 frontiere point, mark size=1.5pt] table[x=k_milieu, y=e_frontiere] {\CodonDonnees/synonyme-energie.csv};
\addlegendentry{\EF{frontier}{frontière}}
\addplot[t3 synonyme, t3 synonyme point, mark size=1.2pt] table[x=k_milieu, y=e_synonyme] {\CodonDonnees/synonyme-energie.csv};
\addlegendentry{\EF{synonym}{synonyme} $\lambda=1$}
\end{axis}
\begin{axis}[t3 graphe, t3 decimales y=1, name=c, at={(a.south west)}, anchor=north west, yshift=-2.25cm,
  width=\linewidth, height=5.6cm, xtick={0,500,...,2500},
  title={(b$'$) \EF{cumulative share of the output error $\sum_{j\le k}\lambda_j e_j$}{part cumulée de l'erreur de sortie $\sum_{j\le k}\lambda_j e_j$}},
  xlabel={\EF{eigendirection $k$}{direction propre $k$}}, ylabel={\EF{share}{part}},
  xmin=0, xmax=2880, ymin=0, ymax=1.02,
  legend style={at={(0.5,-0.3)}, anchor=north}, legend columns=2]
\SynoMuet{0}{1.02}
\addplot[fill=CouleurSynonyme, fill opacity=0.14, draw=none, forget plot] table[x=k, y=part] {\CodonDonnees/synonyme-cumul-aire.csv} \closedcycle;
\addplot[fill=CouleurFrontiere, fill opacity=0.35, draw=none, forget plot] table[x=k, y=part] {\CodonDonnees/synonyme-cumul-aire-frontiere.csv} \closedcycle;
\addplot[CouleurSynonyme, syno calib] table[x=k, y=synonyme_calib] {\CodonDonnees/synonyme-cumul.csv};
\addlegendentry{\EF{synonym, calibration}{synonyme, calibration}}
\addplot[CouleurSynonyme, syno eval] table[x=k, y=synonyme_eval] {\CodonDonnees/synonyme-cumul.csv};
\addlegendentry{\EF{synonym, seen by evaluation}{synonyme, vue par l'évaluation}}
\addplot[CouleurFrontiere, syno calib, line width=1.3pt] table[x=k, y=frontiere_calib] {\CodonDonnees/synonyme-cumul.csv};
\addlegendentry{\EF{frontier, calibration}{frontière, calibration}}
\addplot[CouleurFrontiere, syno eval] table[x=k, y=frontiere_eval] {\CodonDonnees/synonyme-cumul.csv};
\addlegendentry{\EF{frontier, seen by evaluation}{frontière, vue par l'évaluation}}
\draw[t3 rappel] (axis cs:1300,0.66) -- (axis cs:1560,0.36);
\node[t3 nombre=CouleurSynonyme!85!black, anchor=north west, align=left] at (axis cs:1420,0.36)
  {\EF{synonym: area}{synonyme : aire} \CodonNum{\SynoNAireS}\Pv\ \EF{cost seen by evaluation}{coût vu par l'évaluation}
   $\times\,\CodonNum{\SynoNSurS}$};
\draw[t3 rappel] (axis cs:2400,0.965) -- (axis cs:2400,0.2);
\node[t3 nombre=CouleurFrontiereMarque, anchor=north west, align=left] at (axis cs:1420,0.2)
  {\EF{frontier: area}{frontière : aire} \CodonNum{\SynoNAireF}\Pv\ \EF{cost seen by evaluation}{coût vu par l'évaluation}
   $\times\,\CodonNum{\SynoNSurF}$};
\end{axis}
\end{tikzpicture}}
\Titre{SynoBasePropre}{The synonymous mutation in the eigenbasis of $H$}{La mutation synonyme dans la base propre de $H$}
\Legende{SynoBasePropre}{\mes\,\calc}
  {The synonym lowers the calibration cost by pushing the error towards the directions calibration does not see;
   evaluation still sees them.}
  {Le synonyme baisse le coût de calibration en poussant l'erreur vers les directions que la calibration ne voit
   pas ; l'évaluation, elle, les voit.}
  {Before (frontier) and after (synonym $\lambda=1$), for one matrix: layer 8, expert 0, gate of gpt-oss-20b;
   $H=\mathrm{E}[g^2xx^\top]/\mathrm{E}[g^2]$ from the calibration inputs; the output cost is
   $\mathrm{tr}(\Delta W H\Delta W^\top)=\sum_k\lambda_k e_k$. (a) The calibration $H$ has rank
   \CodonNum{\SynoNRang} (the number of calibration inputs): the
   \CodonNum{\the\numexpr\SynoNDirections-\SynoNRang\relax} directions beyond (hatched), out of
   \CodonNum{\SynoNDirections}, weigh nothing in calibration, but the evaluation inputs still weigh them (at least
   \CodonNum{\SynoNLambdaEvalMin}). The coder itself minimised the cost with $H+\delta^*\,\overline{\mathrm{diag}\,H}\,I$,
   $\delta^*=1$ (07 §~2): \CodonNum{\SynoNAmorti} added to every direction, more than the evaluation weight of the
   silent directions; only \SynoNAmortiDessus\ eigenvalues exceed it. (b) The synonym raises the weight error,
   $\sum e_k$ \CodonNum{\SynoNEnergieF} $\to$ \CodonNum{\SynoNEnergieS} ($+\Pourcent{\SynoNEnergieHausse}$), and
   moves it towards the silent directions (\Pourcent{\SynoNPartNulleF} $\to$ \Pourcent{\SynoNPartNulleS} of
   $\sum e_k$ beyond the rank). (b$'$) The cumulative share of the output error, in calibration (solid) and
   seen by the evaluation inputs (dotted). The output cost falls by $\times\,1/\CodonNum{\SynoNCalibRapport}$ in
   calibration (\CodonNum{\SynoNCalibF} $\to$ \CodonNum{\SynoNCalibS}) but only by
   $\times\,1/\CodonNum{\SynoNEvalRapport}$ seen by the evaluation (\CodonNum{\SynoNEvalF} $\to$ \CodonNum{\SynoNEvalS}):
   the frontier's cost seen by the evaluation stays at its calibration cost ($\times\,\CodonNum{\SynoNSurF}$), the
   synonym's rises by $\times\,\CodonNum{\SynoNSurS}$; that is the residual over-fitting. The shaded area between
   the two curves is not its measure: the frontier has one too (\CodonNum{\SynoNAireF}, against
   \CodonNum{\SynoNAireS} for the synonym), because the calibration share reaches 1 at the rank of $H$, for any
   variant. One matrix of one expert; the other two spectra in the same folder are not drawn.}
  {Avant (frontière) et après (synonyme $\lambda=1$), pour une matrice : couche 8, expert 0, gate de gpt-oss-20b ;
   $H=\mathrm{E}[g^2xx^\top]/\mathrm{E}[g^2]$ sur les entrées de calibration ; le coût de sortie vaut
   $\mathrm{tr}(\Delta W H\Delta W^\top)=\sum_k\lambda_k e_k$. (a) Le $H$ de calibration est de rang
   \CodonNum{\SynoNRang} (le nombre d'entrées de calibration) : les
   \CodonNum{\the\numexpr\SynoNDirections-\SynoNRang\relax} directions au-delà (hachurées), sur
   \CodonNum{\SynoNDirections}, ne pèsent rien en calibration, mais les entrées d'évaluation les pèsent encore (au
   moins \CodonNum{\SynoNLambdaEvalMin}). Le codeur, lui, a minimisé le coût avec
   $H+\delta^*\,\overline{\mathrm{diag}\,H}\,I$, $\delta^*=1$ (07 §~2) : \CodonNum{\SynoNAmorti} ajouté à chaque
   direction, plus que le poids d'évaluation des directions muettes ; seules \SynoNAmortiDessus\ valeurs propres le
   dépassent. (b) Le synonyme augmente l'erreur sur les poids, $\sum e_k$ \CodonNum{\SynoNEnergieF} $\to$
   \CodonNum{\SynoNEnergieS} ($+\Pourcent{\SynoNEnergieHausse}$), et la déplace vers les directions muettes
   (\Pourcent{\SynoNPartNulleF} $\to$ \Pourcent{\SynoNPartNulleS} de $\sum e_k$ au-delà du rang).
   (b$'$) La part cumulée de l'erreur de sortie, en calibration (trait plein) et vue par les entrées
   d'évaluation (pointillés). Le coût de sortie baisse de $\times\,1/\CodonNum{\SynoNCalibRapport}$ en calibration
   (\CodonNum{\SynoNCalibF} $\to$ \CodonNum{\SynoNCalibS}) mais seulement de $\times\,1/\CodonNum{\SynoNEvalRapport}$
   vu par l'évaluation (\CodonNum{\SynoNEvalF} $\to$ \CodonNum{\SynoNEvalS}) : le coût de la frontière vu par
   l'évaluation reste à son coût de calibration ($\times\,\CodonNum{\SynoNSurF}$), celui du synonyme monte de
   $\times\,\CodonNum{\SynoNSurS}$ ; c'est le sur-apprentissage résiduel. L'aire ombrée entre les deux courbes n'en
   est pas la mesure : la frontière en a une aussi (\CodonNum{\SynoNAireF}, contre \CodonNum{\SynoNAireS} pour le
   synonyme), parce que la part en calibration atteint 1 au rang de $H$, pour toute variante. Une matrice d'un
   expert ; les deux autres spectres du même dossier ne sont pas tracés.}
\Source{SynoBasePropre}{\fichier{mesures/genome/construction/synonyme/figures/spectre-L08-E00-gate.csv}\Pv{}
  \EF{script}{script} \fichier{memoire/outils/tome3_synonyme.py}.}

% =============================================================================================================
% III-9.2. La queue de la KL : survie sur les 2 550 positions de R1, 768 experts substitués.
% Lignes de quantiles du CSV : 0 Q2_0, 1 frontière, 2 λ = 1, 3 λ = 0 (marques de la forme de leur série).
\newcommand{\SynoQuantile}[2]{% #1 rang dans synonyme-quantiles.csv, #2 style de marque
  \addplot[only marks, #2, forget plot, restrict expr to domain={\coordindex}{#1:#1}]
    table[x=mediane, y=p_mediane] {\CodonDonnees/synonyme-quantiles.csv};
  \addplot[only marks, #2, forget plot, restrict expr to domain={\coordindex}{#1:#1}]
    table[x=p999, y=p_p999] {\CodonDonnees/synonyme-quantiles.csv};}
\pgfplotstableread{\CodonDonnees/synonyme-quantiles.csv}\SynoQuantTable
\newcommand{\GrapheSynoQueue}{%
\pgfplotstablegetelem{1}{mediane}\of\SynoQuantTable\let\SynoMedF\pgfplotsretval
\pgfplotstablegetelem{2}{mediane}\of\SynoQuantTable\let\SynoMedU\pgfplotsretval
\begin{tikzpicture}
\begin{axis}[t3 graphe, width=0.86\linewidth, height=7.4cm, xmode=log, ymode=log,
  xmin=1e-2, xmax=10, ymin=3e-4, ymax=1.4, unbounded coords=jump,
  xlabel={\EF{threshold $t$ (KL per position, nats)}{seuil $t$ (KL par position, nats)}},
  ylabel={$P(\mathrm{KL}>t)$}, legend style={at={(0.03,0.25)}, anchor=south west}]
\draw[black!40, densely dotted, line width=0.5pt] (axis cs:1e-2,0.5) -- (axis cs:10,0.5);
\draw[black!40, densely dotted, line width=0.5pt] (axis cs:1e-2,0.001) -- (axis cs:10,0.001);
\node[t3 etiquette, text=black!65, anchor=south west] at (axis cs:1.1e-2,0.5) {\EF{median}{médiane}};
\node[t3 etiquette, text=black!65, anchor=south west] at (axis cs:1.1e-2,0.001) {\EF{99.9th percentile}{99,9\textsuperscript{e} centile}};
\addplot[t3 qdeux] table[x=t, y=q2_0] {\CodonDonnees/synonyme-survie.csv};
\addlegendentry{llama.cpp Q2\_0}
\addplot[t3 frontiere] table[x=t, y=frontiere] {\CodonDonnees/synonyme-survie.csv};
\addlegendentry{\EF{frontier}{frontière}}
\addplot[t3 synonyme] table[x=t, y=lambda1] {\CodonDonnees/synonyme-survie.csv};
\addlegendentry{\EF{synonym}{synonyme} $\lambda=1$}
\addplot[t3 lambda zero] table[x=t, y=lambda0] {\CodonDonnees/synonyme-survie.csv};
\addlegendentry{\EF{synonym}{synonyme} $\lambda=0$}
\SynoQuantile{0}{t3 qdeux point}
\SynoQuantile{1}{t3 frontiere point}
\SynoQuantile{2}{t3 synonyme point}
\SynoQuantile{3}{t3 lambda zero point}
% le rappel « ÷ » : de la médiane de la frontière à celle de λ = 1, sur la ligne des médianes
\draw[t3 rappel, {Stealth[length=1.6mm]}-, line width=0.6pt, draw=Encre] (axis cs:\SynoMedU,0.32) -- (axis cs:\SynoMedF,0.32);
\draw[t3 rappel] (axis cs:\SynoMedU,0.32) -- (axis cs:\SynoMedU,0.45);
\draw[t3 rappel] (axis cs:\SynoMedF,0.32) -- (axis cs:\SynoMedF,0.45);
\node[t3 etiquette, anchor=north] at (axis cs:0.1,0.31) {$\div\,\CodonNum{\SynoNDivUn}$};
\end{axis}
\end{tikzpicture}}
\Titre{SynoQueue}{The tail of the KL}{La queue de la KL}
\Legende{SynoQueue}{\mes}
  {The synonym divides the median KL by nearly three; the tail rests on too few positions to rank the variants.}
  {Le synonyme divise la KL médiane par près de trois ; la queue repose sur trop peu de positions pour classer les
   variantes.}
  {Survival function $P(\mathrm{KL}>t)$ over the \CodonNum{\SynoNPositions} scored positions of R1 (10 windows), all
   768 experts of gpt-oss-20b replaced in the f64 reference; on the dotted lines, the medians and the 99.9th
   percentiles, each with the mark of its series. The median of the frontier (\CodonNum{\SynoNMedF}) is divided by
   \CodonNum{\SynoNDivUn} with $\lambda=1$ (\CodonNum{\SynoNMedUn}) and by \CodonNum{\SynoNDivZero} with $\lambda=0$
   (\CodonNum{\SynoNMedZero}). Each 99.9th percentile leaves only \SynoNQueueDessus\ positions above it: against
   the frontier, $\lambda=1$ lowers it by \Pourcent{\SynoNQueueUnRun} on R1 but by \Pourcent{\SynoNQueueUnNeuf} on
   the judge's ten new windows, $\lambda=0$ by \Pourcent{\SynoNQueueZeroRun} and \Pourcent{\SynoNQueueZeroNeuf} (new
   windows: frontier through the substituted reference, synonyms through the engine). Between $\lambda=0$ and
   $\lambda=1$ the tail is not comparable at this size (barred, JUGEMENT-construction M2). Left of $10^{-2}$: at most
   \Pourcent{\SynoNSurvieGauche} of the positions of a series. Q2\_0 is computed by llama.cpp on the processor,
   another path (a landmark). Positions of one window share their past and are not independent
   (\cref{fig:t3-syno-foret}).}
  {Fonction de survie $P(\mathrm{KL}>t)$ sur les \CodonNum{\SynoNPositions} positions notées de R1 (10 fenêtres), les
   768 experts de gpt-oss-20b substitués dans la référence f64 ; sur les lignes pointillées, les médianes et les
   99,9\textsuperscript{es} centiles, chacun à la marque de sa série. La médiane de la frontière
   (\CodonNum{\SynoNMedF}) est divisée par \CodonNum{\SynoNDivUn} avec $\lambda=1$ (\CodonNum{\SynoNMedUn}) et par
   \CodonNum{\SynoNDivZero} avec $\lambda=0$ (\CodonNum{\SynoNMedZero}). Chaque 99,9\textsuperscript{e} centile ne
   laisse que \SynoNQueueDessus\ positions au-dessus de lui : contre la frontière, $\lambda=1$ l'abaisse de
   \Pourcent{\SynoNQueueUnRun} sur R1 mais de \Pourcent{\SynoNQueueUnNeuf} sur les dix fenêtres neuves du juge,
   $\lambda=0$ de \Pourcent{\SynoNQueueZeroRun} et de \Pourcent{\SynoNQueueZeroNeuf} (fenêtres neuves : frontière par
   la référence substituée, synonymes par le moteur). Entre $\lambda=0$ et $\lambda=1$, la queue n'est pas
   comparable à cette taille (barré, JUGEMENT-construction M2). À gauche de $10^{-2}$ : au plus
   \Pourcent{\SynoNSurvieGauche} des positions d'une série. Q2\_0 est calculé par llama.cpp sur le processeur, un
   autre chemin (un repère). Les positions d'une fenêtre partagent leur passé et ne sont pas indépendantes
   (\cref{fig:t3-syno-foret}).}
\Source{SynoQueue}{\fichier{mesures/genome/construction/synonyme/kl/kl-*-768-positions.csv},
  \fichier{repere/kl-q2\_0-positions.csv}, \fichier{kl/kl-*-768-juge.json}\Pv{}
  \fichier{mesures/genome/construction/juge/kl/kl-neuf-*-positions.csv}\Pv\ \EF{script}{script}
  \fichier{memoire/outils/tome3_synonyme.py}.}

% =============================================================================================================
% III-9.3 (F-S3 + F-J3). B2 fenêtre par fenêtre : forêt des différences appariées, R1 et fenêtres neuves.
% Un seul axe pour les trois panneaux (bornes écrites par le script, qui s'arrête si un point ou une bande sort).
\newcommand{\SynoForet}[3]{% #1 fichier, #2 titre, #3 1 si les neuves mêlent deux chemins (marques creuses)
\nextgroupplot[title={#2}]
\pgfplotstableread{\CodonDonnees/synonyme-foret-#1-moyennes.csv}\SynoMoy
\pgfplotstablegetelem{0}{moyenne}\of\SynoMoy\let\SynoMa\pgfplotsretval
\pgfplotstablegetelem{0}{deux_sigma}\of\SynoMoy\let\SynoSa\pgfplotsretval
\pgfplotstablegetelem{1}{moyenne}\of\SynoMoy\let\SynoMb\pgfplotsretval
\pgfplotstablegetelem{1}{deux_sigma}\of\SynoMoy\let\SynoSb\pgfplotsretval
\edef\SynoBandes{%
  \noexpand\fill[CouleurRun!14] (axis cs:\SynoMa-\SynoSa,13.6) rectangle (axis cs:\SynoMa+\SynoSa,23.4);
  \noexpand\draw[CouleurRun, line width=1.1pt] (axis cs:\SynoMa,13.6) -- (axis cs:\SynoMa,23.4);
  \noexpand\fill[CouleurNeuves!16] (axis cs:\SynoMb-\SynoSb,0.6) rectangle (axis cs:\SynoMb+\SynoSb,10.4);
  \noexpand\draw[CouleurNeuves, line width=1.1pt] (axis cs:\SynoMb,0.6) -- (axis cs:\SynoMb,10.4);}
\SynoBandes
\draw[black!55, dashed, line width=0.6pt] (axis cs:0,0) -- (axis cs:0,25);
\addplot[only marks, t3 run point, unbounded coords=discard]
  table[x=d, y expr={\thisrow{y}>11 ? \thisrow{y} : nan}] {\CodonDonnees/synonyme-foret-#1.csv};
\ifnum#3=1
\addplot[only marks, mark=diamond, mark size=2.4pt, mark options={solid, draw=CouleurNeuves!85!black, line width=0.7pt},
  unbounded coords=discard] table[x=d, y expr={\thisrow{y}<11 ? \thisrow{y} : nan}] {\CodonDonnees/synonyme-foret-#1.csv};
\else
\addplot[only marks, t3 neuves point, unbounded coords=discard]
  table[x=d, y expr={\thisrow{y}<11 ? \thisrow{y} : nan}] {\CodonDonnees/synonyme-foret-#1.csv};
\fi
}
\newcommand{\GrapheSynoForet}{%
\begin{tikzpicture}
\begin{groupplot}[group style={group size=3 by 1, horizontal sep=0.35cm, y descriptions at=edge left},
  t3 graphe, t3 decimales x=1, width=0.37\linewidth, height=9.2cm, ymin=0, ymax=25, scaled ticks=false,
  xmin=\SynoNForXmin, xmax=\SynoNForXmax, xtick={-0.3,-0.2,-0.1,0},
  ytick={1,...,10,14,15,...,23}, yticklabels={9,8,7,6,5,4,3,2,1,0,9,8,7,6,5,4,3,2,1,0},
  yticklabel style={font=\scriptsize}, ymajorgrids=false, xmajorgrids=true,
  xlabel={$\bar d$ \EF{per window}{par fenêtre}}, ylabel={\EF{window}{fenêtre}}]
\SynoForet{un-frontiere}{$\lambda=1$ \EF{vs}{contre} \EF{frontier}{frontière}}{1}
\node[t3 nombre=CouleurRun, anchor=north west] at (axis cs:\SynoNForXmin+0.01,24.8) {R1\Dp\ \CodonNum{\SynoNForUnFRunZ}\,$\sigma$, \SynoNForUnFRunNeg/10 $<0$};
\node[t3 nombre=CouleurNeuves!85!black, anchor=west, align=left] at (axis cs:\SynoNForXmin+0.01,12)
  {\EF{new}{neuves}\Dp\ \CodonNum{\SynoNForUnFNeufZ}\,$\sigma$, \SynoNForUnFNeufNeg/10 $<0$\\[-1pt]\EF{two paths, indicative}{deux chemins, indicatif}};
\SynoForet{zero-un}{$\lambda=0$ \EF{vs}{contre} $\lambda=1$}{0}
\node[t3 nombre=CouleurRun, anchor=north west] at (axis cs:\SynoNForXmin+0.01,24.8) {R1\Dp\ \CodonNum{\SynoNForZeroUnRunZ}\,$\sigma$, \SynoNForZeroUnRunNeg/10 $<0$};
\node[t3 nombre=CouleurNeuves!85!black, anchor=west] at (axis cs:\SynoNForXmin+0.01,12) {\EF{new}{neuves}\Dp\ \CodonNum{\SynoNForZeroUnNeufZ}\,$\sigma$, \SynoNForZeroUnNeufNeg/10 $<0$};
\SynoForet{zero-frontiere}{$\lambda=0$ \EF{vs}{contre} \EF{frontier}{frontière}}{0}
\node[t3 nombre=CouleurRun, anchor=north west] at (axis cs:\SynoNForXmin+0.01,24.8) {R1\Dp\ \CodonNum{\SynoNForZeroFRunZ}\,$\sigma$, \SynoNForZeroFRunNeg/10 $<0$};
\node[t3 nombre=CouleurNeuves!85!black, anchor=west] at (axis cs:\SynoNForXmin+0.01,12) {\EF{new}{neuves}\Dp\ \CodonNum{\SynoNForZeroFNeufZ}\,$\sigma$, \SynoNForZeroFNeufNeg/10 $<0$};
\end{groupplot}
\node[rotate=90, anchor=south, text=CouleurRun, font=\sffamily\small\bfseries] at ($(group c1r1.west)+(-1.0cm,2.0cm)$) {R1};
\node[rotate=90, anchor=south, text=CouleurNeuves!85!black, font=\sffamily\small\bfseries] at ($(group c1r1.west)+(-1.0cm,-2.1cm)$) {\EF{new}{neuves}};
\end{tikzpicture}}
\Titre{SynoForet}{Criterion B2, window by window}{Le critère B2, fenêtre par fenêtre}
\Legende{SynoForet}{\mes}
  {Both synonyms beat the frontier on both sets of windows, each in at least \SynoNForBatMin\ of
   \SynoNForFenetres\ windows; between them the gap is much smaller, and $\lambda=0$ passes on R1 by a hair.}
  {Les deux synonymes battent la frontière sur les deux jeux de fenêtres, chacun dans au moins \SynoNForBatMin\
   fenêtres sur \SynoNForFenetres\ ; entre eux, l'écart est bien plus petit, et $\lambda=0$ ne passe que de
   justesse sur R1.}
  {Paired differences of the mean KL per window, $\bar d=\overline{\mathrm{KL}}_b-\overline{\mathrm{KL}}_a$, negative
   when the second variant wins; one axis for the three panels, no point outside the frame. Top (circles), the ten
   windows of R1 by lot B through the substituted f64 reference (for $\lambda=0$ against the frontier, the judge's
   recomputation on the same path). Bottom (diamonds), ten new windows of the judge: panel 1 sets $\lambda=1$
   through the engine against the frontier through the substituted reference, two paths, published as indicative
   only (JUGEMENT-construction §~9; hollow diamonds; on R1, where both paths exist, this mix moves $\bar d$ by
   \CodonNum{\SynoNForMelange}); panel 2 through the engine; panel 3 through the substituted reference. Vertical
   line and band: mean and $\pm2\sigma$, $\sigma$ the standard deviation of the ten windows divided by $\sqrt{10}$.
   $\lambda=0$ against $\lambda=1$ passes on R1 by a hair (\CodonNum{\SynoNForZeroUnRunZ}\,$\sigma$, window 7
   against) and with more margin on the new windows (\CodonNum{\SynoNForZeroUnNeufZ}\,$\sigma$). The new windows are
   drawn from wiki.test, the corpus of R1: they test other text, not another kind of text.}
  {Différences appariées des KL moyennes par fenêtre, $\bar d=\overline{\mathrm{KL}}_b-\overline{\mathrm{KL}}_a$,
   négatives quand la seconde variante gagne ; un seul axe pour les trois panneaux, aucun point hors du cadre. En
   haut (cercles), les dix fenêtres de R1 par le lot B, par la référence f64 substituée (pour $\lambda=0$ contre la
   frontière, le recalcul du juge par le même chemin). En bas (losanges), dix fenêtres neuves du juge : le
   panneau 1 met $\lambda=1$ par le moteur contre la frontière par la référence substituée, deux chemins, publiés à
   titre indicatif seulement (JUGEMENT-construction §~9 ; losanges creux ; sur R1, où les deux chemins existent, ce
   mélange déplace $\bar d$ de \CodonNum{\SynoNForMelange}) ; le panneau 2 par le moteur ; le panneau 3 par la
   référence substituée. Trait vertical et bande : moyenne et $\pm2\sigma$, $\sigma$ l'écart-type des dix fenêtres
   divisé par $\sqrt{10}$. $\lambda=0$ contre $\lambda=1$ passe de justesse sur R1
   (\CodonNum{\SynoNForZeroUnRunZ}\,$\sigma$, la fenêtre 7 contre) et avec plus de marge sur les fenêtres neuves
   (\CodonNum{\SynoNForZeroUnNeufZ}\,$\sigma$). Les fenêtres neuves sont tirées de wiki.test, le corpus de R1 :
   elles éprouvent un autre texte, pas une autre sorte de texte.}
\Source{SynoForet}{\fichier{mesures/genome/construction/synonyme/kl/apparie-*-par-fenetre.csv},
  \fichier{mesures/genome/construction/juge/figures/apparie-par-fenetre.csv}\Pv\ \EF{two-path check}{mesure du
  mélange} \fichier{mesures/genome/construction/synonyme/moteur/kl-moteur-*-positions.csv}\Pv\ \EF{script}{script}
  \fichier{memoire/outils/tome3_synonyme.py}.}

% =============================================================================================================
% III-9.6 (F-S8). La frontière des variantes à 768 experts (réserve R3 : aucun point du protocole à 32 experts).
% Lignes de synonyme-variantes.csv : 0 frontière, 1 λ = 1, 2 λ = 0, 3 Q2_0, 4 exact. Le cadre (G7 du juge de
% design) ne tient que les trois premières ; Q2_0 et le moteur exact sont des repères de marge.
\newcommand{\SynoVariante}[2]{% #1 rang dans synonyme-variantes.csv, #2 styles (marque, couleur des barres)
  \addplot[only marks, restrict expr to domain={\coordindex}{#1:#1}, erreurs, /pgfplots/.cd, #2]
    table[x=bits, y=kl, y error=deux_sigma_f] {\CodonDonnees/synonyme-variantes.csv};}
\newcommand{\GrapheSynoVariantes}{%
\begin{tikzpicture}
\begin{axis}[t3 graphe, t3 points, t3 decimales x=2, width=0.8\linewidth, height=7.2cm, ymode=log, clip=false,
  xmin=\SynoNVarXmin, xmax=\SynoNVarXmax, ymin=\SynoNVarYmin, ymax=\SynoNVarYmax,
  ytick={0.1,0.15,0.2,0.3}, log ticks with fixed point, yticklabel style={/pgf/number format/.cd, fixed, fixed zerofill, precision=2},
  xlabel={\EF{bits per expert weight (real file)}{bits par poids d'expert (fichier réel)}},
  ylabel={\EF{end-to-end KL (768 experts)}{KL de bout en bout (768 experts)}}]
\SynoVariante{0}{t3 frontiere point, error bars/error bar style={CouleurFrontiereMarque, line width=0.6pt}}
\SynoVariante{1}{t3 synonyme point, error bars/error bar style={CouleurSynonyme, line width=0.6pt}}
\SynoVariante{2}{t3 lambda zero point, error bars/error bar style={CouleurLambdaZero, line width=0.6pt}}
\draw[-{Stealth[length=2.2mm]}, Encre, line width=0.7pt] (axis cs:\SynoNVarFBits-0.006,\SynoNVarFKL*0.97) -- (axis cs:\SynoNVarUnBits+0.008,\SynoNVarUnKL*1.06);
\node[t3 etiquette, anchor=south, rotate=0] at (axis cs:0.5*\SynoNVarFBits+0.5*\SynoNVarUnBits,0.205) {\EF{dominates on both axes}{domine sur les deux axes}};
\node[t3 nombre=CouleurFrontiereMarque, anchor=west] at (axis cs:\SynoNVarFBits+0.006,\SynoNVarFKL) {\EF{frontier}{frontière}};
% Troisième passe (G7 du juge de design) : le cadre resserré (2,56 à 2,70 bits, tiré des données par le script)
% ne laisse plus de place à gauche de λ = 1 : son étiquette passe à droite de son point ; celle de λ = 0 passe sous
% son point, à gauche de sa barre.
\node[t3 nombre=CouleurSynonyme, anchor=west] at (axis cs:\SynoNVarUnBits+0.004,\SynoNVarUnKL) {\EF{synonym}{synonyme} $\lambda=1$};
\node[t3 nombre=CouleurLambdaZero!85!black, anchor=north east, align=right] at (axis cs:\SynoNVarZeroBits-0.004,\SynoNVarZeroKL*0.975)
  {\EF{synonym}{synonyme} $\lambda=0$\\$-\CodonNum{\SynoNAchatKL}$ \EF{KL for}{de KL pour} $+\CodonNum{\SynoNAchatBits}$ bit};
% repères de marge : une flèche au bord, vers le point hors du cadre
% (le repère est dessiné hors du cadre, au-delà du coin, pour ne pas se lire comme un point du cadre)
\draw[-{Stealth[length=2.2mm]}, CouleurQdeux!85!black, line width=0.7pt] (rel axis cs:0.07,0.9) -- (rel axis cs:-0.015,1.02);
\path[t3 qdeux point] plot coordinates {(rel axis cs:-0.03,1.045)};
\node[t3 nombre=CouleurQdeux!80!black, anchor=north west, align=left] at (rel axis cs:0.075,0.94)
  {llama.cpp Q2\_0, \EF{off the frame}{hors du cadre}\\\num[round-mode=places, round-precision=2]{\SynoNVarQBits} bits,
   KL \num[round-mode=places, round-precision=2]{\SynoNVarQKL}};
% Le moteur exact : son étiquette en bas, à gauche de la barre de λ = 0 (N1 du juge des chiffres de la construction :
% elle en couvrait le bas) ; la flèche part de son bord droit, sous la barre.
\draw[-{Stealth[length=2.2mm]}, Encre, line width=0.7pt] (rel axis cs:0.775,0.055) -- (rel axis cs:1.02,0.015);
\path[t3 exact point] plot coordinates {(rel axis cs:1.045,0.0)};
\node[t3 nombre=Encre, anchor=south east, align=right] at (rel axis cs:0.765,0.025)
  {\EF{exact engine (v1.0), off the frame}{moteur exact (v1.0), hors du cadre}\\\CodonNum{\SynoNExactBits} bits, KL \CodonNum{\SynoNExactKL}};
\end{axis}
\end{tikzpicture}}
\Titre{SynoVariantes}{The frontier of the variants, 768 experts}{La frontière des variantes, 768 experts}
\Legende{SynoVariantes}{\mes}
  {Synonym $\lambda=1$ dominates the frontier on both axes; $\lambda=0$ buys KL with bits.}
  {Le synonyme $\lambda=1$ domine la frontière sur les deux axes ; $\lambda=0$ achète de la KL avec des bits.}
  {All 768 experts of gpt-oss-20b replaced: bits per expert weight of the real companion file (useful bits)
   against the end-to-end KL on the ten windows of R1, logarithmic; bars: $\pm2\sigma$ by windows, recomputed from
   the per-position files (for the frontier $\sigma=\CodonNum{\SynoNVarFSf}$, against \CodonNum{\SynoNVarFSn} by
   positions: the per-position $\sigma$ published in 07 is \CodonNum{\SynoNSigmaRapportMin} to
   \CodonNum{\SynoNSigmaRapportMax} times too small). $\lambda=0$ buys \CodonNum{\SynoNAchatKL} of KL with
   \CodonNum{\SynoNAchatBits} bit. Margin landmarks, outside the frame (arrows): llama.cpp Q2\_0 (its KL computed
   by llama.cpp on the processor, another path than the substituted reference) and the exact engine against the
   f64 reference (useful bits of the v1.0 file, the definition used for the companions). Reservation R3: these
   points are measured with 768 experts replaced; the points of the round (\cref{t3syn:ch}) are measured with 32 experts,
   and nothing joins them to these.}
  {Les 768 experts de gpt-oss-20b substitués : bits par poids d'expert du fichier compagnon réel (bits utiles)
   contre la KL de bout en bout sur les dix fenêtres de R1, en échelle logarithmique ; barres : $\pm2\sigma$ par
   fenêtres, recalculé depuis les fichiers par position (pour la frontière $\sigma=\CodonNum{\SynoNVarFSf}$, contre
   \CodonNum{\SynoNVarFSn} par positions : le $\sigma$ par positions publié dans 07 est de
   \CodonNum{\SynoNSigmaRapportMin} à \CodonNum{\SynoNSigmaRapportMax} fois trop petit). $\lambda=0$ achète
   \CodonNum{\SynoNAchatKL} de KL avec \CodonNum{\SynoNAchatBits} bit. Repères de marge, hors du cadre (flèches) :
   llama.cpp Q2\_0 (sa KL calculée par llama.cpp sur le processeur, un autre chemin que la référence substituée) et
   le moteur exact contre la référence f64 (bits utiles du fichier v1.0, la définition prise pour les compagnons).
   Réserve R3 : ces points sont mesurés avec les 768 experts substitués ; ceux de la ronde (\cref{t3syn:ch}) le sont avec 32
   experts, et rien ne les relie à ceux-ci.}
\Source{SynoVariantes}{\fichier{mesures/genome/construction/synonyme/figures/variantes-768.csv},
  \fichier{kl/kl-*-768-positions.csv}, \fichier{repere/kl-q2\_0-positions.csv}\Pv{}
  \fichier{mesures/moteur/O/journaux/o3/c4-20b-comparer.json}\Pv\ \fichier{mesures/genome/construction/juge/J-T0-fichiers.json}\Pv\ \EF{script}{script} \fichier{memoire/outils/tome3_synonyme.py}.}
\fi
